\subsection{Recovering boundary sewing from the completed tower}

\begin{result}{E.35}{Proved}{Intrinsic framed sewing and detection of the open-group actions}
The labelled completed exit system with its clutching,
forgetting and composition maps recovers the full completed
boundary-fixed mapping class groupoids and their boundary
sewing maps.  The recovery is functorial for equivalences
of that system and coherent for iterated gluing and
self-gluing.  Its groupoid comparison with the marked-pants
operad of E.27 and E.34 is restriction to the bordered
cores of the labelled tripod construction.

Consequently an equivalence of this completed system whose
genus-zero restriction is the identity has trivial induced
outer action on every stable open group
$\widehat\Gamma_{g,n}$.  This conclusion uses neither a
congruence-completion identification nor additional covering
maps.  The reconstruction of all noninvertible exit arrows
from this induced group action remains the separate
crossing-map question.
\end{result}
\begin{proof}
We describe the sewing map with all of its auxiliary
projections.  Encode each parametrized boundary of a
bordered core by an attached rational tripod, with two
distinct external labels, as in E.29 and E.33.
For the annular unit use its single-node four-point-sphere
crossing instead: the two annular boundary twists both
map to that crossing's positive meridian.  Attaching two
tripods to an annulus contracts its two-valent core and
gives this same single node.  Thus the unit is represented
without treating its two interface curves as distinct.
Sewing this annulus is the usual identity collar extension.
The following construction concerns stable bordered cores.
Choose one boundary on each of two cores, and clutch
one external marked point of each of their tripods.
Call their original core-interface nodes $i,j$ and
the new middle node $m$; the other external marked
point on each tripod is its companion.
The same construction applies to two boundaries of
one core.  Retain the other boundary-tag tripods.
All branches are pinned and the distinct companion
labels distinguish the two middle tripods.

Let $P$ be the product of the boundary-fixed core
groups in the two-core case, and the single such
group in the self-gluing case.  Cutting every node
of this graph gives the literal presentation
\[
                 L=P\times\Z_m.                             \tag{E.120}
\]
Indeed each tripod contributes its three commuting
boundary twists.  A core-interface gluing identifies
its inner twist with the corresponding core twist;
the middle gluing identifies the two middle twists;
ordinary marked-point capping kills the companion
twists.  Eliminating these generators leaves precisely
(E.120).  Conversely extend each core diffeomorphism
by the identity, and realize the remaining generator
by the middle annular twist.  Restriction and these
extensions are inverse on every displayed generator.

E.29 realizes $\widehat L$ as the automorphism group
of the full crossing from this source to the open
stratum.  It can equivalently be read on a full
adjacent smoothing flag: all the labelled nodes are
fixed individually, and E.32 identifies both groups
with the same full completion of the discrete
stabilizer.  Thus changing their smoothing order
does not change that identification.

There are two independently specified projections
\[
 q:\widehat L\longrightarrow\widehat P,
 \qquad r:\widehat L\longrightarrow\widehat\Z_m.              \tag{E.121}
\]
To obtain $q$, order a smoothing flag so that the
middle node is smoothed last and omit that last
arrow.  The remaining flag is exactly the clutching
of the input tripod flags.  Its stabilizer is
$P=L/\Z_m$, and E.32 gives its full completion.
The identifications with the input factors are
also the actual clutching maps; their full topology
can be seen on each target component from E.33.

To obtain $r$, smooth the middle node first and
retain every other node at that first step.
Normalize at the two core-interface nodes and
select the connected component formed by the
two middle tripods.  It is a labelled four-point
sphere with its one-node crossing.  The normalized
core-side branches are ordinary marked points;
their twists are therefore killed, whereas the
middle twist is retained.  The resulting map is
$r(P)=1$, $r(\tau_m)=\tau_m$.

This is a map supplied by the retained system,
rather than an extra choice of a normal-coordinate
section.  At the first adjacent crossing, the
clutching description of the normalized target
identifies its link group with the product of
the unchanged core groups and the four-point
link group.  Projecting that product is exactly
the indicated component-selection normalization.
The component labels are fixed.  The comparison
uses the actual clutching homomorphisms on its
source and target, and E.32 identifies the full
completed link.  Thus an equivalence preserving
clutching and the pinned composition data also
preserves this projection.  Equivalently it is
the admissible normalization of A.9 and A.15
applied before the two core-interface nodes are
smoothed; those nodes are retained throughout
that square.

The product map $(q,r)$ is the completed isomorphism
(E.120).  Hence there is a unique homomorphism
\[
 s:\widehat P\longrightarrow\widehat L,
              \qquad qs=\id,\quad rs=1,                       \tag{E.122}
\]
and it is the extension-by-identity section on the
middle annulus.  This characterizes the section
using the two retained maps.  In particular a
possible additional central character on $P$
cannot be inserted into this section: its value
would be detected by $r$.

Forget the two companion labels.  Their three-node
chain stabilizes to one sewn node $e$.  In local
plumbing coordinates the actual stabilization map
is
\[
       q_e=q_iq_mq_j,
       \qquad (a_i,a_m,a_j)\longmapsto a_i+a_m+a_j
              \quad\hbox{on meridian lattices}.               \tag{E.123}
\]
One obtains the formula by substituting the three
equations $xy=q$ successively along the two
unstable rational components; their normalized
coordinates are fixed by their three special
points before forgetting.  It is also two
successive instances of C.1's bridge formula.
The map on crossing groupoids is the completed
forgetful functor of E.1 and E.4.  Its projection
to the bordered target group, obtained by
smoothing $e$ while retaining the other boundary
tags, gives after composition with $s$
\[
 \widehat\Gamma_{g,r}^{\partial}\times
 \widehat\Gamma_{h,t}^{\partial}
       \longrightarrow\widehat\Gamma_{g+h,r+t-2}^{\partial},
                                                               \tag{E.124}
\]
or the corresponding self-gluing map.  Formula
(E.123) shows that both chosen boundary twists
map to the same positive sewn twist.  The
middle coordinate is zero by (E.122).
On every core diffeomorphism the map is ordinary
extension over the sewn collars.  These statements
identify (E.124), on every element, with the
usual boundary-sewing homomorphism.

For several simultaneous sewings the construction
has
$L=P\times\Z^{E_{\mathrm{mid}}}$, with one
four-point projection $r_e$ for each middle node.
The section is uniquely characterized by $q s=\id$
and $r_e s=1$ for every $e$.  Each order of
successive constructions gives this same section.
Stabilization multiplies the same normal parameters,
or adds the same integer meridians, in each order.
Disjoint normalizations and the permitted component
projections commute by A.15.  These facts prove
associativity, compatibility of self-gluings, and
relabeling coherence, with the stated comparison
maps.  All equations remain equations after full
completion by E.1.

We now make the groupoid and kernel conclusion
explicit.  For each marked pants graph $G$ take
its bordered surface $S_G$, attach the labelled
tripods, and take its distinguished full crossing
$\eta_G$.  Add these finitely many representatives
to the chosen skeleton in each type.  The full
groupoid of commutative squares between these
crossings has morphisms exactly the bordered
mapping classes $S_G\to S_H$: restrict a square
to its core, and inversely extend a core map by
the standard tripods.  Its endpoint on the
boundary stratum is the capped core map, and
the full square retains all boundary twists.
These two operations are inverse.  E.33 proves
that the same full completion, with its entire
morphism torsor, occurs inside the completed
arrow groupoid.  Thus this is the groupoid
comparison with $\widehat{\mathcal S}$, not only
an identification of automorphism groups.
The sewing just constructed is its ordinary
sewing by (E.124).

An equivalence of the tower induces an equivalence
of these groupoids with their sewing comparison
maps.  For an element of the genus-zero kernel,
choose its coherent genus-zero identification
with the identity.  The object operad of the
marked-pants model is free on the two cyclically
marked pants objects, with the cylinder as unit,
as in \cite[Definition 3.1]{ET_BonattoRobertson2026}.
Extend those chosen identifications to each
graph object by simultaneously sewing all its vertex pants,
using the operation comparison maps.  The chosen genus-zero
identification respects the $C_3$ stabilizer of each cyclic
pants generator.  The
coherences just proved make the resulting
identification independent of the order of
sewing and of vertex and edge relabeling, and the cyclic identifications on the
generating pants are already fixed in genus
zero.  Conjugating by these object identifications
gives a strict object-fixing automorphism of
$\widehat{\mathcal S}$ whose genus-zero restriction
is the identity.  E.27 makes it the identity.
In particular the original equivalence induces
an inner automorphism on every completed
boundary-fixed core group.

For $n>0$, forgetting one companion from each
of the $n$ tripods realizes the surjection
\[
 \widehat\Gamma_{g,n}^{\partial}
          \twoheadrightarrow\widehat\Gamma_{g,n}.             \tag{E.125}
\]
The induced maps commute with that actual
forgetful functor.  An inner action on its
source therefore gives an inner action on
its target.  For stable closed types, first
use type $(g,1)$ and then forget that mark;
the completed Birman quotient is onto, so
the same argument applies.  This proves
the stated all-type kernel assertion.
It establishes the induced open-group action;
no inference about a remaining arrow displacement
is made from this quotient argument.
\end{proof}
